% Options for packages loaded elsewhere
\PassOptionsToPackage{unicode}{hyperref}
\PassOptionsToPackage{hyphens}{url}
\PassOptionsToPackage{dvipsnames,svgnames,x11names}{xcolor}
%
\documentclass[
  letterpaper,
  DIV=11,
  numbers=noendperiod]{scrartcl}

\usepackage{amsmath,amssymb}
\usepackage{iftex}
\ifPDFTeX
  \usepackage[T1]{fontenc}
  \usepackage[utf8]{inputenc}
  \usepackage{textcomp} % provide euro and other symbols
\else % if luatex or xetex
  \usepackage{unicode-math}
  \defaultfontfeatures{Scale=MatchLowercase}
  \defaultfontfeatures[\rmfamily]{Ligatures=TeX,Scale=1}
\fi
\usepackage{lmodern}
\ifPDFTeX\else  
    % xetex/luatex font selection
\fi
% Use upquote if available, for straight quotes in verbatim environments
\IfFileExists{upquote.sty}{\usepackage{upquote}}{}
\IfFileExists{microtype.sty}{% use microtype if available
  \usepackage[]{microtype}
  \UseMicrotypeSet[protrusion]{basicmath} % disable protrusion for tt fonts
}{}
\makeatletter
\@ifundefined{KOMAClassName}{% if non-KOMA class
  \IfFileExists{parskip.sty}{%
    \usepackage{parskip}
  }{% else
    \setlength{\parindent}{0pt}
    \setlength{\parskip}{6pt plus 2pt minus 1pt}}
}{% if KOMA class
  \KOMAoptions{parskip=half}}
\makeatother
\usepackage{xcolor}
\setlength{\emergencystretch}{3em} % prevent overfull lines
\setcounter{secnumdepth}{-\maxdimen} % remove section numbering
% Make \paragraph and \subparagraph free-standing
\makeatletter
\ifx\paragraph\undefined\else
  \let\oldparagraph\paragraph
  \renewcommand{\paragraph}{
    \@ifstar
      \xxxParagraphStar
      \xxxParagraphNoStar
  }
  \newcommand{\xxxParagraphStar}[1]{\oldparagraph*{#1}\mbox{}}
  \newcommand{\xxxParagraphNoStar}[1]{\oldparagraph{#1}\mbox{}}
\fi
\ifx\subparagraph\undefined\else
  \let\oldsubparagraph\subparagraph
  \renewcommand{\subparagraph}{
    \@ifstar
      \xxxSubParagraphStar
      \xxxSubParagraphNoStar
  }
  \newcommand{\xxxSubParagraphStar}[1]{\oldsubparagraph*{#1}\mbox{}}
  \newcommand{\xxxSubParagraphNoStar}[1]{\oldsubparagraph{#1}\mbox{}}
\fi
\makeatother

\usepackage{color}
\usepackage{fancyvrb}
\newcommand{\VerbBar}{|}
\newcommand{\VERB}{\Verb[commandchars=\\\{\}]}
\DefineVerbatimEnvironment{Highlighting}{Verbatim}{commandchars=\\\{\}}
% Add ',fontsize=\small' for more characters per line
\usepackage{framed}
\definecolor{shadecolor}{RGB}{241,243,245}
\newenvironment{Shaded}{\begin{snugshade}}{\end{snugshade}}
\newcommand{\AlertTok}[1]{\textcolor[rgb]{0.68,0.00,0.00}{#1}}
\newcommand{\AnnotationTok}[1]{\textcolor[rgb]{0.37,0.37,0.37}{#1}}
\newcommand{\AttributeTok}[1]{\textcolor[rgb]{0.40,0.45,0.13}{#1}}
\newcommand{\BaseNTok}[1]{\textcolor[rgb]{0.68,0.00,0.00}{#1}}
\newcommand{\BuiltInTok}[1]{\textcolor[rgb]{0.00,0.23,0.31}{#1}}
\newcommand{\CharTok}[1]{\textcolor[rgb]{0.13,0.47,0.30}{#1}}
\newcommand{\CommentTok}[1]{\textcolor[rgb]{0.37,0.37,0.37}{#1}}
\newcommand{\CommentVarTok}[1]{\textcolor[rgb]{0.37,0.37,0.37}{\textit{#1}}}
\newcommand{\ConstantTok}[1]{\textcolor[rgb]{0.56,0.35,0.01}{#1}}
\newcommand{\ControlFlowTok}[1]{\textcolor[rgb]{0.00,0.23,0.31}{\textbf{#1}}}
\newcommand{\DataTypeTok}[1]{\textcolor[rgb]{0.68,0.00,0.00}{#1}}
\newcommand{\DecValTok}[1]{\textcolor[rgb]{0.68,0.00,0.00}{#1}}
\newcommand{\DocumentationTok}[1]{\textcolor[rgb]{0.37,0.37,0.37}{\textit{#1}}}
\newcommand{\ErrorTok}[1]{\textcolor[rgb]{0.68,0.00,0.00}{#1}}
\newcommand{\ExtensionTok}[1]{\textcolor[rgb]{0.00,0.23,0.31}{#1}}
\newcommand{\FloatTok}[1]{\textcolor[rgb]{0.68,0.00,0.00}{#1}}
\newcommand{\FunctionTok}[1]{\textcolor[rgb]{0.28,0.35,0.67}{#1}}
\newcommand{\ImportTok}[1]{\textcolor[rgb]{0.00,0.46,0.62}{#1}}
\newcommand{\InformationTok}[1]{\textcolor[rgb]{0.37,0.37,0.37}{#1}}
\newcommand{\KeywordTok}[1]{\textcolor[rgb]{0.00,0.23,0.31}{\textbf{#1}}}
\newcommand{\NormalTok}[1]{\textcolor[rgb]{0.00,0.23,0.31}{#1}}
\newcommand{\OperatorTok}[1]{\textcolor[rgb]{0.37,0.37,0.37}{#1}}
\newcommand{\OtherTok}[1]{\textcolor[rgb]{0.00,0.23,0.31}{#1}}
\newcommand{\PreprocessorTok}[1]{\textcolor[rgb]{0.68,0.00,0.00}{#1}}
\newcommand{\RegionMarkerTok}[1]{\textcolor[rgb]{0.00,0.23,0.31}{#1}}
\newcommand{\SpecialCharTok}[1]{\textcolor[rgb]{0.37,0.37,0.37}{#1}}
\newcommand{\SpecialStringTok}[1]{\textcolor[rgb]{0.13,0.47,0.30}{#1}}
\newcommand{\StringTok}[1]{\textcolor[rgb]{0.13,0.47,0.30}{#1}}
\newcommand{\VariableTok}[1]{\textcolor[rgb]{0.07,0.07,0.07}{#1}}
\newcommand{\VerbatimStringTok}[1]{\textcolor[rgb]{0.13,0.47,0.30}{#1}}
\newcommand{\WarningTok}[1]{\textcolor[rgb]{0.37,0.37,0.37}{\textit{#1}}}

\providecommand{\tightlist}{%
  \setlength{\itemsep}{0pt}\setlength{\parskip}{0pt}}\usepackage{longtable,booktabs,array}
\usepackage{calc} % for calculating minipage widths
% Correct order of tables after \paragraph or \subparagraph
\usepackage{etoolbox}
\makeatletter
\patchcmd\longtable{\par}{\if@noskipsec\mbox{}\fi\par}{}{}
\makeatother
% Allow footnotes in longtable head/foot
\IfFileExists{footnotehyper.sty}{\usepackage{footnotehyper}}{\usepackage{footnote}}
\makesavenoteenv{longtable}
\usepackage{graphicx}
\makeatletter
\newsavebox\pandoc@box
\newcommand*\pandocbounded[1]{% scales image to fit in text height/width
  \sbox\pandoc@box{#1}%
  \Gscale@div\@tempa{\textheight}{\dimexpr\ht\pandoc@box+\dp\pandoc@box\relax}%
  \Gscale@div\@tempb{\linewidth}{\wd\pandoc@box}%
  \ifdim\@tempb\p@<\@tempa\p@\let\@tempa\@tempb\fi% select the smaller of both
  \ifdim\@tempa\p@<\p@\scalebox{\@tempa}{\usebox\pandoc@box}%
  \else\usebox{\pandoc@box}%
  \fi%
}
% Set default figure placement to htbp
\def\fps@figure{htbp}
\makeatother
% definitions for citeproc citations
\NewDocumentCommand\citeproctext{}{}
\NewDocumentCommand\citeproc{mm}{%
  \begingroup\def\citeproctext{#2}\cite{#1}\endgroup}
\makeatletter
 % allow citations to break across lines
 \let\@cite@ofmt\@firstofone
 % avoid brackets around text for \cite:
 \def\@biblabel#1{}
 \def\@cite#1#2{{#1\if@tempswa , #2\fi}}
\makeatother
\newlength{\cslhangindent}
\setlength{\cslhangindent}{1.5em}
\newlength{\csllabelwidth}
\setlength{\csllabelwidth}{3em}
\newenvironment{CSLReferences}[2] % #1 hanging-indent, #2 entry-spacing
 {\begin{list}{}{%
  \setlength{\itemindent}{0pt}
  \setlength{\leftmargin}{0pt}
  \setlength{\parsep}{0pt}
  % turn on hanging indent if param 1 is 1
  \ifodd #1
   \setlength{\leftmargin}{\cslhangindent}
   \setlength{\itemindent}{-1\cslhangindent}
  \fi
  % set entry spacing
  \setlength{\itemsep}{#2\baselineskip}}}
 {\end{list}}
\usepackage{calc}
\newcommand{\CSLBlock}[1]{\hfill\break\parbox[t]{\linewidth}{\strut\ignorespaces#1\strut}}
\newcommand{\CSLLeftMargin}[1]{\parbox[t]{\csllabelwidth}{\strut#1\strut}}
\newcommand{\CSLRightInline}[1]{\parbox[t]{\linewidth - \csllabelwidth}{\strut#1\strut}}
\newcommand{\CSLIndent}[1]{\hspace{\cslhangindent}#1}

\KOMAoption{captions}{tableheading}
\makeatletter
\@ifpackageloaded{caption}{}{\usepackage{caption}}
\AtBeginDocument{%
\ifdefined\contentsname
  \renewcommand*\contentsname{Table of contents}
\else
  \newcommand\contentsname{Table of contents}
\fi
\ifdefined\listfigurename
  \renewcommand*\listfigurename{List of Figures}
\else
  \newcommand\listfigurename{List of Figures}
\fi
\ifdefined\listtablename
  \renewcommand*\listtablename{List of Tables}
\else
  \newcommand\listtablename{List of Tables}
\fi
\ifdefined\figurename
  \renewcommand*\figurename{Figure}
\else
  \newcommand\figurename{Figure}
\fi
\ifdefined\tablename
  \renewcommand*\tablename{Table}
\else
  \newcommand\tablename{Table}
\fi
}
\@ifpackageloaded{float}{}{\usepackage{float}}
\floatstyle{ruled}
\@ifundefined{c@chapter}{\newfloat{codelisting}{h}{lop}}{\newfloat{codelisting}{h}{lop}[chapter]}
\floatname{codelisting}{Listing}
\newcommand*\listoflistings{\listof{codelisting}{List of Listings}}
\makeatother
\makeatletter
\makeatother
\makeatletter
\@ifpackageloaded{caption}{}{\usepackage{caption}}
\@ifpackageloaded{subcaption}{}{\usepackage{subcaption}}
\makeatother

\usepackage{bookmark}

\IfFileExists{xurl.sty}{\usepackage{xurl}}{} % add URL line breaks if available
\urlstyle{same} % disable monospaced font for URLs
\hypersetup{
  pdftitle={Outcome-Blind Estimator Validation for Targeted Learning in a Clean-Room Pharmacoepidemiology Simulation},
  pdfauthor={Andrew N. Mertens},
  colorlinks=true,
  linkcolor={blue},
  filecolor={Maroon},
  citecolor={Blue},
  urlcolor={Blue},
  pdfcreator={LaTeX via pandoc}}


\title{Outcome-Blind Estimator Validation for Targeted Learning in a
Clean-Room Pharmacoepidemiology Simulation}
\author{Andrew N. Mertens}
\date{Invalid Date}

\begin{document}
\maketitle


\section{1. Introduction}\label{introduction}

\subsection{1.1 Motivation}\label{motivation}

Real-world data (RWD) studies in pharmacoepidemiology increasingly
support regulatory safety evaluations, yet the credibility of these
analyses depends on the transparency and rigor of the analytic workflow.
In conventional observational studies, analytic decisions---variable
selection, model specification, positivity handling, sensitivity
analyses---are often made iteratively after inspecting outcome data,
creating researcher degrees of freedom that can inflate false-positive
rates and undermine reproducibility {[}@Lash2021; @Hernan2016{]}.

Clean-room study designs address this problem by imposing governance
structures that separate analytic roles, restrict outcome access, and
require pre-specification of the analysis plan before outcome data are
examined. This approach draws on principles from registered reports and
pre-analysis plans in clinical trials, adapted for the complexity of
observational data analysis {[}@Nosek2018; @Berger2021{]}.

This paper describes a simulation study that evaluates the performance
of targeted learning estimators---specifically TMLE and AIPW with Super
Learner---within a staged clean-room framework. The simulation is
calibrated to a realistic pharmacoepidemiologic safety question: whether
sofosbuvir-containing (SOF) hepatitis C virus (HCV) treatment regimens
are associated with elevated risk of acute kidney injury (AKI) compared
with non-SOF direct-acting antiviral (DAA) regimens.

\subsection{1.2 The clean-room staged
framework}\label{the-clean-room-staged-framework}

We organize the analysis into four stages:

\begin{itemize}
\tightlist
\item
  \textbf{Stage 1 (Pre-specification)}: Define the causal question,
  target trial, estimand, identification assumptions, and estimation
  strategy. Document all analytic decisions in a pre-registered
  Statistical Analysis Plan (SAP).
\item
  \textbf{Stage 2 (Outcome-blind diagnostics)}: With treatment and
  covariate data available but outcomes withheld, evaluate propensity
  score overlap, effective sample size (ESS), standardized mean
  differences (SMD), and weight behavior. Confirm that the planned
  analysis is feasible without modifying the estimand based on outcome
  information.
\item
  \textbf{Stage 3 (Outcome-blind simulation)}: Using a pre-specified
  data-generating process (DGP) that reflects the target study's
  covariate structure and plausible treatment-outcome relationships, run
  the full estimation pipeline on simulated data. Evaluate bias,
  coverage, and variance calibration of the planned estimators.
  Validate---rather than discover---the analytic plan's adequacy.
\item
  \textbf{Stage 4 (Unblinded estimation)}: Apply the validated pipeline
  to the real data with outcome access. Report results alongside
  pre-specified diagnostics and sensitivity analyses.
\end{itemize}

\subsubsection{1.2.1 Conceptual components: governance, outcome-blind
diagnostics, outcome-blind
simulation}\label{conceptual-components-governance-outcome-blind-diagnostics-outcome-blind-simulation}

Readers should distinguish four components that serve different roles in
this framework:

\textbf{Clean-room governance} refers to the organizational and
access-control structure: role separation between data providers,
methodologists, quality assurance, privacy compliance, and reporting
teams. Governance rules determine who has access to which data at each
stage, require documentation of analytic decisions, and enforce sign-off
procedures before results are released. Clean-room governance is a
workflow framework, not a statistical method.

\textbf{Outcome-blind diagnostics} are analyses performed during Stage 2
that evaluate the feasibility of the planned analysis without access to
outcome data. These include propensity score overlap assessment,
effective sample size under inverse probability weights, standardized
mean differences before and after weighting, and examination of weight
distributions. Outcome-blind diagnostics can identify practical
positivity violations, covariate balance problems, and data quality
issues that would compromise any downstream estimation strategy.

\textbf{Outcome-blind simulation} is the Stage 3 activity in which the
full estimation pipeline is applied to synthetic data generated from a
pre-specified DGP. The purpose is not to discover the best estimator
post hoc, but to validate that the pre-specified analysis
plan---including estimator choice, Super Learner library, truncation
parameters, and variance estimation---performs adequately under
conditions resembling the target study. Simulation results that reveal
poor coverage or substantial bias trigger revision of the analysis plan
\emph{before} outcomes are accessed, preserving the integrity of
pre-specification.

\textbf{Targeted learning estimation} (TMLE, AIPW) is the statistical
framework applied at Stage 4 to produce causal effect estimates. The
targeting step, efficient influence function, and double robustness
properties are statistical properties of the estimator, not features of
the governance framework.

These four components are complementary but conceptually distinct.
Clean-room governance structures the workflow; outcome-blind diagnostics
and simulations validate the analytic plan; targeted learning provides
the estimation machinery.

\section{2. Methods}\label{methods}

\subsection{2.1 Target trial framing and treatment
strategy}\label{target-trial-framing-and-treatment-strategy}

Causal inference from observational data proceeds by emulating a
hypothetical randomized trial---the \textbf{target trial}
{[}@Hernan2016{]}. Making the target trial explicit clarifies the
scientific question and disciplines analytic choices.

\textbf{Eligibility.} Adults aged \(\geq\) 18 years with chronic HCV who
are eligible for DAA treatment, with at least 365 days of continuous
enrollment prior to the index date, no prior AKI event, and no prior HCV
treatment (SOF or non-SOF).

\textbf{Treatment strategies.} Two static treatment strategies are
compared: (1) initiate a SOF-containing DAA regimen at baseline; (2)
initiate a non-SOF DAA regimen at baseline.

\textbf{Time zero.} The index date is the date of treatment initiation.
Eligibility criteria and baseline covariates are assessed using the
365-day lookback period.

\textbf{Treatment strategy and switching.} The primary estimand
corresponds to a hypothetical intervention in which individuals initiate
the assigned regimen at baseline and are followed regardless of
subsequent treatment changes. Treatment switching is handled as a
censoring event (see Section 2.5), and the estimand is interpreted as
the effect of initial treatment assignment in a setting where switching
is eliminated---that is, censoring at switch is treated as an
intervention that removes switching from the observed data. This
corresponds to an ``initiator'' or ``intention-to-treat-like'' estimand
under hypothetical elimination of switching, rather than a sustained
treatment strategy. We note this distinction explicitly because
different handling of switching (e.g., per-protocol, as-treated) would
correspond to different estimands.

\textbf{Follow-up.} From time zero until the first AKI event,
administrative censoring, treatment switching (plus a 30-day risk
window), or 720 days, whichever comes first.

\textbf{Outcome.} Cause-specific cumulative incidence of AKI at
pre-specified evaluation times \(t^* \in \{90, 180\}\) days. These time
points were chosen because 90 days approximates one treatment course for
most DAA regimens (capturing acute drug-related toxicity), while 180
days captures any delayed nephrotoxic effects that may emerge after
treatment completion.

\subsection{2.2 Target parameter}\label{target-parameter}

The causal estimand is the \textbf{risk difference} at a pre-specified
time \(t^*\), under hypothetical elimination of censoring:

\[\psi(t^*) = E\bigl[Y_{t^*}^{a=1, \bar{c}=0}\bigr] - E\bigl[Y_{t^*}^{a=0, \bar{c}=0}\bigr]\]

where \(Y_{t^*}^{a, \bar{c}=0}\) denotes the potential AKI indicator at
time \(t^*\) under treatment strategy \(a\) and under a hypothetical
intervention that eliminates all censoring (both administrative and
switching-related). The superscript \(\bar{c} = 0\) represents a
hypothetical world in which no individual is censored before
\(t^*\)---this is an intervention on the censoring process, not a data
manipulation.

This estimand targets \textbf{cause-specific cumulative incidence}.
Competing events (e.g., death from non-AKI causes) are not explicitly
modeled in the current simulation. In the DGP, event times represent AKI
occurrence only, and censoring encompasses administrative loss to
follow-up and treatment switching. The absence of competing risks is a
simplification; in a real-world application, death as a competing event
would require either a composite outcome, a competing-risks framework
(e.g., Fine-Gray subdistribution hazard or cause-specific cumulative
incidence function), or explicit causal modeling of the competing event.
We note this as a limitation of the current simulation (see Discussion).

Optionally, we also report the risk ratio
\(E[Y_{t^*}^{a=1, \bar{c}=0}] / E[Y_{t^*}^{a=0, \bar{c}=0}]\).

\subsection{2.3 Identification
assumptions}\label{identification-assumptions}

Identification of the causal estimand from observed data requires the
following assumptions:

\textbf{Consistency.} The observed outcome for an individual who
received treatment \(a\) equals the counterfactual outcome under
intervention \(a\): if \(A_i = a\), then \(Y_i = Y_i^{a}\). This
requires that the treatment is well-defined and that there are no hidden
treatment versions.

\textbf{Conditional exchangeability (no unmeasured confounding).} The
potential outcomes are independent of treatment assignment, conditional
on measured baseline covariates:

\[Y^{a, \bar{c}=0} \perp\!\!\!\perp A \mid W\]

In our simulation, confounding operates exclusively through baseline
covariates \(W\) (demographics, clinical history, concomitant
medications). There are no time-varying confounders between treatment
initiation and outcome. Therefore, conditional exchangeability must hold
given the baseline covariate vector \(W\). In settings with time-varying
confounders affecting both censoring and outcome, the required
assumption would be \emph{sequential} conditional exchangeability: at
each time \(t\), the potential outcomes are independent of treatment and
censoring conditional on the full observed history through \(t\).

For censoring, we additionally require:

\[Y^{a} \perp\!\!\!\perp C_t \mid A, W \quad \text{(independent censoring given } A, W\text{)}\]

In the dependent-censoring scenario, this assumption is satisfied
because censoring depends only on \(A\) and \(W\) (through the outcome
linear predictor), both of which are conditioned on. In a real-world
analysis, this assumption requires that all factors driving loss to
follow-up and treatment switching are captured in the measured covariate
set.

\textbf{Positivity.} Within every stratum of the measured confounders
that has positive probability in the target population, every treatment
level must have a positive probability of being assigned:

\[0 < P(A = 1 \mid W = w) < 1 \quad \text{for all } w \text{ with } P(W = w) > 0\]

Note that positivity is defined with respect to \emph{strata of measured
confounders}, not with respect to individual-level probabilities in
isolation. Two types of positivity problems arise in practice:

\begin{itemize}
\item
  \textbf{Theoretical (structural) violations} occur when certain
  covariate combinations deterministically preclude a treatment. For
  example, if a specific drug is contraindicated in patients with severe
  renal impairment, no adjustment method can estimate the treatment
  effect in that stratum because the counterfactual is undefined by
  design. These violations represent a fundamental identification
  problem.
\item
  \textbf{Practical (random) violations} occur when certain covariate
  strata could in principle receive either treatment but do so rarely in
  the observed data. These lead to extreme inverse probability weights,
  inflated variance, and unstable estimates. The causal effect is
  formally identified but may be poorly estimable.
\end{itemize}

\textbf{No interference.} One individual's treatment does not affect
another individual's outcome. This is plausible for individual-level
drug treatments.

\subsubsection{2.3.1 Connecting assumptions to Stage-2
diagnostics}\label{connecting-assumptions-to-stage-2-diagnostics}

Several Stage-2 (outcome-blind) diagnostics directly evaluate the
plausibility and practical implications of these assumptions:

\begin{longtable}[]{@{}
  >{\raggedright\arraybackslash}p{(\linewidth - 4\tabcolsep) * \real{0.3333}}
  >{\raggedright\arraybackslash}p{(\linewidth - 4\tabcolsep) * \real{0.3333}}
  >{\raggedright\arraybackslash}p{(\linewidth - 4\tabcolsep) * \real{0.3333}}@{}}
\toprule\noalign{}
\begin{minipage}[b]{\linewidth}\raggedright
Assumption
\end{minipage} & \begin{minipage}[b]{\linewidth}\raggedright
Diagnostic
\end{minipage} & \begin{minipage}[b]{\linewidth}\raggedright
What it reveals
\end{minipage} \\
\midrule\noalign{}
\endhead
\bottomrule\noalign{}
\endlastfoot
Positivity & Propensity score overlap plot & Separation near 0 or 1
indicates near-violations \\
Positivity & Effective sample size (ESS) &
\(n_{\text{eff}} = (\sum w_i)^2 / \sum w_i^2\); low ESS indicates that a
few observations dominate \\
Positivity & Maximum inverse probability weight & Extreme weights signal
instability \\
Exchangeability (partial) & Weighted SMD & Residual imbalance after
weighting suggests incomplete confounding adjustment \\
Positivity + truncation & Truncation sensitivity & Varying truncation
bounds evaluates the bias-variance trade-off \\
\end{longtable}

These diagnostics cannot verify exchangeability (which depends on
unmeasured confounders), but they can identify settings where the
analysis is likely to be unreliable regardless of the estimation method.

\subsection{2.4 Data-generating process}\label{data-generating-process}

The simulation DGP is implemented in \texttt{R/dgp/DGP.R} via the
function \texttt{generate\_hcv\_data()}. We describe each component with
sufficient detail for independent reproduction.

\subsubsection{2.4.1 Baseline covariates
(W)}\label{baseline-covariates-w}

A cohort of \(N = 125{,}000\) individuals is generated with:

\begin{itemize}
\tightlist
\item
  \textbf{Demographics}: age \(\sim \max(N(48, 13), 18)\); sex (58\%
  male); race (5 categories: white 48\%, Black 14\%, Hispanic 6\%, Asian
  2\%, other 30\%); US region (NE 20\%, MW 18\%, S 37\%, W 25\%);
  enrollment duration \(\sim \text{Poisson}(420)\) days.
\item
  \textbf{Clinical history} (independent Bernoulli): CKD (8\%), prior
  AKI (5\%), cirrhosis (18\%), portal hypertension (4\%), ESLD (2\%),
  heart failure (7\%), hypertension (45\%), HIV (4\%), sepsis (3\%),
  diabetes (20\%), cancer (8\%), and others.
\item
  \textbf{Concomitant medications} (independent Bernoulli): NSAIDs
  (25\%), ACE inhibitors/ARBs (30\%), diuretics (22\%), aminoglycosides
  (5\%), contrast agents (8\%), statins (15\%), and others.
\item
  \textbf{Continuous}: BMI \(\sim N(28, 5)\).
\end{itemize}

\textbf{Eligibility exclusions} remove individuals with enrollment
\textless{} 365 days, age \textless{} 18, prior AKI, or prior HCV
treatment (SOF or non-SOF).

\subsubsection{2.4.2 Treatment assignment
model}\label{treatment-assignment-model}

Treatment \(A \in \{0, 1\}\) (SOF vs.~non-SOF) is assigned via a
logistic model with target prevalence \(p_{\text{SOF}} \approx 0.36\).

\textbf{Simple scenario} (\texttt{complexity\ =\ FALSE}):
\[\text{logit}\, P(A = 1 \mid W) = \alpha_0 + 0.015 \cdot \text{age} + 0.30 \cdot \text{cirrhosis} + 0.25 \cdot \text{CKD} + 0.15 \cdot \text{HIV} + 0.10 \cdot \text{diabetes} - 0.10 \cdot \text{cancer} + \epsilon_i\]
where \(\epsilon_i \sim N(0, 0.6)\) and
\(\alpha_0 = \text{logit}(p_{\text{SOF}}) - \overline{\text{lp}}_0\)
ensures the marginal treatment prevalence is approximately
\(p_{\text{SOF}}\).

\textbf{Complex scenario} (\texttt{complexity\ =\ TRUE}): the linear
predictor additionally includes:
\[0.02 \cdot \text{BMI}^2/100 - 0.3 \cdot \sin(0.1 \cdot \text{BMI}) + 0.5 \cdot (\text{age}/50)^3 + 1.5 \cdot \text{CKD} \times \text{cancer} + 0.8 \cdot \text{HIV} \times \log(1 + \text{age})\]

These nonlinear and interaction terms ensure that a correctly specified
parametric logistic model is unavailable to the analyst, requiring
flexible estimation.

Propensity scores are clipped to \([0.05, 0.95]\) to enforce positivity
by construction, although practical near-violations remain possible in
the complex scenario.

\subsubsection{2.4.3 Outcome model (AKI event
times)}\label{outcome-model-aki-event-times}

Individual event times are generated from an exponential (proportional
hazards) or piecewise exponential (non-proportional hazards) model:

\[h_i(t) = h_0 \cdot \exp(\eta_i) \cdot \text{HR}(A_i, t)\]

where the baseline log-hazard linear predictor \(\eta_i\) is:

\textbf{Simple} (\texttt{complexity\ =\ FALSE}):
\[\eta_i = -2.8 + 0.03 \cdot \text{age}_i + 0.7 \cdot \text{CKD}_i + 0.5 \cdot \text{cirrhosis}_i + 0.3 \cdot \text{HF}_i + 0.25 \cdot \text{NSAID}_i + 0.20 \cdot \text{contrast}_i\]

\textbf{Complex} (\texttt{complexity\ =\ TRUE}):
\[\eta_i = -2.8 + 0.03 \cdot \text{age}_i + 0.0005 \cdot \text{age}_i^2 + 0.7 \cdot \text{CKD}_i + 0.5 \cdot \text{cirrhosis}_i + 0.02 \cdot \text{BMI}_i^2/100 - 0.3 \cdot \sin(0.1 \cdot \text{BMI}_i) + 0.4 \cdot \text{HF}_i \cdot \text{ACE/ARB}_i + 0.6 \cdot \text{NSAID}_i \cdot A_i + 0.3 \cdot \text{contrast}_i \cdot \log(1 + \text{age}_i)\]

The treatment hazard ratio function \(\text{HR}(A, t)\) is defined as:

\textbf{Proportional hazards} (\texttt{np\_hazard\ =\ FALSE}):
\(\text{HR}(A=1, t) = 1.25\) for all \(t\). Event times are exponential
with rate \(h_0 \cdot \exp(\eta_i) \cdot 1.25^{A_i}\).

\textbf{Non-proportional hazards} (\texttt{np\_hazard\ =\ TRUE}): A
piecewise exponential model with change-point at \(\tau = 45\) days:

\[\text{HR}(A=1, t) = \begin{cases} 1.25 & \text{if } t \leq \tau \\ 0.70 & \text{if } t > \tau \end{cases}\]

Event times are generated from the piecewise exponential distribution
via inversion: for rates \(r_1 = h_0 \exp(\eta_i) \cdot 1.25\) and
\(r_2 = h_0 \exp(\eta_i) \cdot 0.70\), with
\(U \sim \text{Uniform}(0,1)\):
\[T_i = \begin{cases} -\log(1 - U) / r_{1,i} & \text{if } U \leq 1 - e^{-r_{1,i}\tau} \\ \tau - \log\bigl(\frac{1-U}{e^{-r_{1,i}\tau}}\bigr) / r_{2,i} & \text{otherwise} \end{cases}\]

The non-proportional hazards scenario means that standard Cox regression
(which assumes a constant log-HR) will be misspecified for the treatment
effect.

Default baseline hazard parameters: \(h_0 = 3 \times 10^{-4}\) (simple)
or \(h_0 = 5 \times 10^{-5}\) (complex).

\subsubsection{2.4.4 Censoring model}\label{censoring-model}

\textbf{Administrative censoring} (\texttt{dep\_censor\ =\ FALSE}):
\(C_{\text{admin}} \sim \text{Exp}(\text{rate} = 1/100)\), independent
of covariates and treatment.

\textbf{Dependent (informative) censoring}
(\texttt{dep\_censor\ =\ TRUE}):
\[C_{\text{admin}} \sim \text{Exp}\bigl(\text{rate} = (1/100) \cdot \exp(0.04 \cdot \eta_i + 0.03 \cdot A_i)\bigr)\]

where \(\eta_i\) is the outcome linear predictor. This creates
informative censoring: patients at higher baseline AKI risk and
SOF-treated patients are censored at higher rates. The censoring depends
on \(A\) and \(W\) only through variables included in the adjustment
set, so the conditional independent censoring assumption remains
identifiable given \((A, W)\).

A hard administrative cutoff is applied at \(T_{\max} = 720\) days.

\subsubsection{2.4.5 Treatment switching
model}\label{treatment-switching-model}

Treatment switching is implemented as \textbf{censoring at crossover}
rather than as a modeled time-varying treatment exposure. This means the
estimand approximates the treatment effect under hypothetical
elimination of switching, conditional on the switching mechanism being
correctly specified.

\textbf{Simple switching} (\texttt{switch\_on\ =\ FALSE}): Treatment
duration \(\sim \text{Poisson}(84)\) for SOF, \(\text{Poisson}(70)\) for
non-SOF. Switching is rare (3\% Bernoulli) and induces censoring at
switch time plus a 30-day risk window.

\textbf{Informative switching} (\texttt{switch\_on\ =\ TRUE}): Switching
time follows a survival process:
\[\text{Switch time} \sim \text{Exp}(\lambda_{\text{sw}})\]
\[\lambda_{\text{sw}} = \lambda_0 \cdot \exp(\gamma_A \cdot A_i + \gamma_{\text{CKD}} \cdot \text{CKD}_i)\]

with \(\lambda_0 = 2.5 \times 10^{-5}\), \(\gamma_A = 0.60\),
\(\gamma_{\text{CKD}} = 0.40\). Patients are censored at (switch time +
30-day risk window) or \(T_{\max}\), whichever is earlier. The switching
mechanism depends on both treatment assignment and CKD status, creating
informative censoring that requires adjustment through the censoring
model \(\hat{g}_C\).

\subsubsection{2.4.6 Observed data
structure}\label{observed-data-structure}

For each individual, the observed follow-up time is:
\[\tilde{T}_i = \min(T_i^{\text{event}},\; C_i^{\text{admin}},\; C_i^{\text{switch}})\]

The event indicator is:
\[\Delta_i = I(T_i^{\text{event}} \leq \tilde{T}_i)\]

where \(I(\cdot)\) is the indicator function, equal to 1 if the AKI
event occurred before any censoring event, and 0 otherwise. Latent
variables (true event times, censoring times) are excluded from the
analysis dataset.

\subsection{2.5 Estimators}\label{estimators}

\subsubsection{2.5.1 TMLE (targeted maximum likelihood
estimation)}\label{tmle-targeted-maximum-likelihood-estimation}

TMLE is a \textbf{substitution estimator} that produces causal effect
estimates by evaluating an updated outcome regression at the
counterfactual treatment values {[}@vanderLaan2006; @vanderLaan2011{]}.
The procedure has three stages:

\begin{enumerate}
\def\labelenumi{\arabic{enumi}.}
\item
  \textbf{Initial estimation of nuisance parameters.} Two nuisance
  functions are estimated:

  \begin{itemize}
  \tightlist
  \item
    The outcome regression: \(\bar{Q}(A, W) = E[Y_{t^*} \mid A, W]\)
  \item
    The propensity score (treatment mechanism):
    \(g(W) = P(A = 1 \mid W)\)
  \item
    The censoring mechanism: \(g_C(A, W) = P(C = 0 \mid A, W)\) (in
    survival settings)
  \end{itemize}

  These are estimated using Super Learner (see Section 2.6).
\item
  \textbf{Targeting step.} The initial outcome regression
  \(\hat{\bar{Q}}\) is updated along a least-favorable parametric
  submodel so that the empirical mean of the \textbf{efficient influence
  function (EIF)} is approximately zero. The fluctuation submodel for
  bounded outcomes is:
  \[\text{logit}\,\hat{\bar{Q}}^*(A, W) = \text{logit}\,\hat{\bar{Q}}(A, W) + \epsilon \cdot H(A, W)\]
  where \(H(A, W) = A/\hat{g}(W) - (1-A)/(1-\hat{g}(W))\) is the clever
  covariate, and \(\epsilon\) is estimated by maximum likelihood. The
  targeting step is what distinguishes TMLE from a naive plug-in
  estimator: it adjusts the initial fit to reduce bias for the specific
  target parameter, ensuring that the estimating equation corresponding
  to the EIF is solved.
\item
  \textbf{Substitution.} The ATE estimate is obtained by substituting
  the \emph{updated} regression into the parameter mapping:
  \[\hat{\psi}_{\text{TMLE}} = \frac{1}{n}\sum_{i=1}^n \bigl[\hat{\bar{Q}}^*(1, W_i) - \hat{\bar{Q}}^*(0, W_i)\bigr]\]
\end{enumerate}

Because TMLE solves the EIF equation, it achieves:

\begin{itemize}
\tightlist
\item
  \textbf{Double robustness}: consistency if either \(\hat{\bar{Q}}\) or
  \(\hat{g}\) is consistently estimated.
\item
  \textbf{Semiparametric efficiency}: when both nuisance models are
  consistently estimated, TMLE achieves the smallest possible asymptotic
  variance among regular asymptotically linear estimators.
\item
  \textbf{Asymptotic linearity}:
  \(\hat{\psi} - \psi_0 = n^{-1}\sum_{i} \text{EIF}(O_i) + o_p(n^{-1/2})\),
  enabling Wald-type inference.
\end{itemize}

As a substitution estimator, TMLE guarantees that estimates remain
within the parameter space (e.g., risk estimates between 0 and 1).

\subsubsection{2.5.2 TMLE with cross-fitting
(TMLE-CF)}\label{tmle-with-cross-fitting-tmle-cf}

Cross-fitting (also called sample splitting) decouples nuisance
parameter training from evaluation, mitigating Donsker-class and
empirical process conditions that are required for theoretical
guarantees when highly adaptive learners are used {[}@Chernozhukov2018;
@Zheng2011{]}. Data are partitioned into \(V\) folds; for each fold,
nuisance parameters are estimated on the remaining \(V-1\) folds and
predictions are generated on the held-out fold. The targeting step and
substitution are then performed on the full-sample predictions. This is
the approach used in \textbf{cross-validated TMLE (CV-TMLE)}
{[}@Zheng2011{]}.

Cross-fitting is particularly important when the Super Learner library
includes highly adaptive algorithms (e.g., HAL, deep learners) that may
not satisfy the entropy conditions required by standard TMLE theory.

\subsubsection{2.5.3 AIPW (augmented inverse probability
weighting)}\label{aipw-augmented-inverse-probability-weighting}

AIPW is a doubly robust estimator that, like TMLE, combines outcome
regression and propensity score estimation. However, AIPW is \textbf{not
a substitution estimator}: it directly evaluates the EIF rather than
updating an outcome model. The AIPW estimate is:

\[\hat{\psi}_{\text{AIPW}} = \frac{1}{n}\sum_{i=1}^n \Biggl[\frac{A_i(Y_i - \hat{\bar{Q}}(1, W_i))}{\hat{g}(W_i)} - \frac{(1-A_i)(Y_i - \hat{\bar{Q}}(0, W_i))}{1 - \hat{g}(W_i)} + \hat{\bar{Q}}(1, W_i) - \hat{\bar{Q}}(0, W_i)\Biggr]\]

AIPW shares the double-robustness and semiparametric efficiency
properties of TMLE. A practical limitation is that, because AIPW is not
a substitution estimator, its estimates may \textbf{fall outside the
parameter space} (e.g., estimated risk differences outside \([-1, 1]\),
or estimated risks outside \([0, 1]\)). This can occur when propensity
scores are extreme or the sample is small.

\subsubsection{2.5.4 Cox proportional hazards
regression}\label{cox-proportional-hazards-regression}

Cox PH regression estimates a \textbf{conditional hazard ratio}
comparing SOF to non-SOF:

\[h(t \mid A, W) = h_0(t) \exp(\beta A + W\gamma)\]

where \(\beta\) is the (assumed constant) log-hazard ratio for
treatment. Cox PH does \textbf{not} directly estimate the target
parameter \(\psi(t^*)\) (a marginal risk difference at a specific time
\(t^*\)). It estimates a different quantity---a conditional hazard
ratio---under a proportional hazards assumption. It is included as a
comparator to represent the standard pharmacoepidemiologic approach and
to demonstrate the consequences of target parameter misalignment and
model misspecification (particularly when the true treatment effect is
non-proportional). Results from Cox PH should be interpreted with this
caveat.

\subsubsection{2.5.5 Unadjusted comparison}\label{unadjusted-comparison}

An unadjusted risk difference (crude comparison of AKI rates between SOF
and non-SOF groups) serves as a reference to quantify the magnitude of
confounding bias.

\subsection{2.6 Super Learner for nuisance parameter
estimation}\label{super-learner-for-nuisance-parameter-estimation}

Super Learner is an ensemble machine learning method that selects the
optimal weighted combination of candidate algorithms by minimizing
\textbf{cross-validated risk} {[}@vanderLaan2007; @Polley2010{]}.

The analyst specifies:

\begin{itemize}
\tightlist
\item
  A \textbf{loss function} appropriate to the prediction task:

  \begin{itemize}
  \tightlist
  \item
    \emph{Squared error loss} for continuous outcomes
  \item
    \emph{Negative log-likelihood (Bernoulli deviance)} for binary
    outcomes (used for \(\bar{Q}\))
  \item
    \emph{Logistic loss} for the treatment model \(g(W)\)
  \end{itemize}
\item
  A \textbf{library of candidate learners}, including both parametric
  and non-parametric algorithms:

  \begin{itemize}
  \tightlist
  \item
    GLM (generalized linear model)
  \item
    GAM (generalized additive model, via \texttt{mgcv})
  \item
    Random forest (via \texttt{ranger})
  \item
    XGBoost (gradient boosted trees)
  \item
    BART (Bayesian additive regression trees, via \texttt{dbarts})
  \item
    HAL (highly adaptive lasso, optional)
  \item
    \texttt{SL.mean} (intercept-only model, as a baseline reference)
  \end{itemize}
\end{itemize}

Super Learner finds the convex combination of candidate learners that
minimizes the cross-validated risk. By including a diverse library, the
ensemble adapts to the complexity of the true nuisance functions without
requiring the analyst to select a single model specification.

\subsection{2.7 Outcome-blind simulation
phase}\label{outcome-blind-simulation-phase}

The outcome-blind simulation phase (Stage 3) serves as a \textbf{design
validation and pre-specification tool}. Its purpose is to confirm that
the pre-specified analysis plan---estimator choice, Super Learner
library, truncation bounds, variance estimation, and diagnostic
criteria---performs adequately under realistic data-generating
conditions before outcome data are accessed.

The outcome-blind simulation phase does not ``discover'' the best
estimator. Rather, it validates the feasibility and robustness of
pre-specified analytic choices. If simulations reveal that the planned
estimator exhibits poor coverage or substantial bias under plausible DGP
scenarios, the analysis plan is revised \emph{before} unblinding. This
revision is principled and documented, preserving the integrity of
pre-specification.

Key validation targets include:

\begin{itemize}
\tightlist
\item
  \textbf{Bias}: Does the estimator recover the known true treatment
  effect (Monte Carlo ground truth)?
\item
  \textbf{Coverage}: Do 95\% confidence intervals contain the true
  parameter value at approximately the nominal rate?
\item
  \textbf{Variance calibration}: Do influence-curve-based standard error
  estimates agree with the empirical standard deviation across
  simulation replicates?
\item
  \textbf{Robustness}: Is performance maintained across DGP scenarios of
  varying complexity (proportional vs.~non-proportional hazards,
  independent vs.~dependent censoring, linear vs.~non-linear nuisance
  functions)?
\end{itemize}

The simulation validates the full pipeline, not just the point
estimator---variance estimation, truncation, and diagnostics are all
exercised.

\section{3. Simulation Design}\label{simulation-design}

\subsection{3.1 Design}\label{design}

The simulation evaluates five DGP scenarios that progressively introduce
realistic complications:

\begin{longtable}[]{@{}
  >{\raggedright\arraybackslash}p{(\linewidth - 8\tabcolsep) * \real{0.2000}}
  >{\raggedright\arraybackslash}p{(\linewidth - 8\tabcolsep) * \real{0.2000}}
  >{\raggedright\arraybackslash}p{(\linewidth - 8\tabcolsep) * \real{0.2000}}
  >{\raggedright\arraybackslash}p{(\linewidth - 8\tabcolsep) * \real{0.2000}}
  >{\raggedright\arraybackslash}p{(\linewidth - 8\tabcolsep) * \real{0.2000}}@{}}
\toprule\noalign{}
\begin{minipage}[b]{\linewidth}\raggedright
Scenario
\end{minipage} & \begin{minipage}[b]{\linewidth}\raggedright
Confounding
\end{minipage} & \begin{minipage}[b]{\linewidth}\raggedright
Censoring
\end{minipage} & \begin{minipage}[b]{\linewidth}\raggedright
Switching
\end{minipage} & \begin{minipage}[b]{\linewidth}\raggedright
Outcome model
\end{minipage} \\
\midrule\noalign{}
\endhead
\bottomrule\noalign{}
\endlastfoot
1. Unconfounded & None (randomized) & Independent & None & Simple, PH \\
2. Baseline confounding & Measured baseline & Independent & None &
Simple, PH \\
3. Informative censoring & Measured baseline & Dependent on \((A, W)\) &
None & Complex, PH \\
4. Treatment switching & Measured baseline & Dependent on \((A, W)\) &
Informative (depends on \(A\), CKD) & Complex, non-PH \\
5. Positivity stress & Measured baseline, strong & Dependent on
\((A, W)\) & Informative & Complex, non-PH \\
\end{longtable}

Each scenario is run for \textbf{200 simulation replicates} at
\(n = 50{,}000\) per replicate. The choice of 200 replicates reflects a
trade-off between computational cost (each replicate involves Super
Learner with cross-validation) and Monte Carlo precision. With 200
replicates, the Monte Carlo standard error for a coverage estimate of
0.95 is approximately \(\sqrt{0.95 \times 0.05 / 200} \approx 0.015\),
meaning coverage estimates have a 95\% Monte Carlo confidence interval
of roughly \(\pm 3\) percentage points. This is sufficient to detect
substantial departures from nominal coverage (e.g., 88\% vs.~95\%) but
may not resolve small differences (e.g., 93\% vs.~95\%). The
implications for interpreting coverage results are discussed in Section
4.

\subsubsection{3.1.1 Ground truth}\label{ground-truth}

The Monte Carlo ``true'' treatment effect for each scenario is computed
by generating \(N = 200{,}000\) observations under each counterfactual
treatment assignment (\texttt{treat\_override\ =\ "all\_treated"} and
\texttt{treat\_override\ =\ "all\_control"}) and computing the
difference in mean outcomes at each \(t^*\). With \(N = 200{,}000\), the
Monte Carlo standard error of the ground truth is on the order of
\(\sqrt{p(1-p)/200{,}000}\), which for typical AKI rates
(\(p \approx 0.01\text{--}0.05\)) is approximately
\(2\text{--}5 \times 10^{-4}\). This Monte Carlo error is small relative
to the estimator standard errors (which are on the order of
\(10^{-2}\text{--}10^{-3}\)), but it is not zero. Any discrepancy
between the Monte Carlo truth and the true parameter value will appear
as a small systematic component in bias calculations.

\subsection{3.2 Evaluation metrics}\label{evaluation-metrics}

For each scenario, we compute:

\begin{itemize}
\tightlist
\item
  \textbf{Bias}: \(\overline{\hat{\psi}} - \psi_{\text{true}}\), where
  \(\overline{\hat{\psi}}\) is the average estimate across replicates
  and \(\psi_{\text{true}}\) is the Monte Carlo ground truth.
\item
  \textbf{Root mean squared error (RMSE)}:
  \(\sqrt{n_{\text{rep}}^{-1} \sum_r (\hat{\psi}_r - \psi_{\text{true}})^2}\).
\item
  \textbf{Empirical SE vs.~IC-based SE}: Comparison of the standard
  deviation of \(\hat{\psi}\) across replicates with the average
  estimated standard error (from the EIF).
\item
  \textbf{95\% CI coverage}: Proportion of replicates in which the
  Wald-type 95\% CI contains \(\psi_{\text{true}}\).
\end{itemize}

\section{4. Results}\label{results}

\emph{{[}Simulation results, tables, and figures to be populated after
running the simulation pipeline. The following subsections describe how
results should be interpreted.{]}}

\subsection{4.1 Interpreting coverage and Monte Carlo
error}\label{interpreting-coverage-and-monte-carlo-error}

Coverage is computed as the fraction of simulation replicates in which
the estimated 95\% confidence interval contains the Monte Carlo ground
truth \(\psi_{\text{true}}\). Several factors complicate interpretation:

\textbf{Monte Carlo variability in coverage estimates.} With 200
replicates, coverage is itself a random variable. A nominal 95\%
coverage rate will fluctuate: the 95\% Monte Carlo confidence interval
for a true coverage of 0.95 is approximately \([0.92, 0.98]\). Coverage
estimates within this range are consistent with nominal performance and
should not be over-interpreted. Coverage below 0.90 or above 0.99 is
more informative.

\textbf{Near-positivity violations.} When propensity scores are close to
0 or 1, the clever covariate \(H(A, W)\) becomes large, inflating both
the point estimate variance and the EIF-based variance estimate. In
finite samples, the EIF variance estimator may underestimate the true
variance, leading to under-coverage. Truncation mitigates this at the
cost of introducing bias.

\textbf{Truncation effects.} Truncating propensity scores at
\([0.01, 0.99]\) introduces a bias-variance trade-off. Coverage can
deteriorate if the truncation-induced bias exceeds the reduction in
variance, particularly in the positivity-stress scenario.

\textbf{EIF variance estimation.} The influence-curve-based variance
estimator assumes that the nuisance parameter estimates converge at
rates sufficient for the remainder term to be negligible. When flexible
learners are used without cross-fitting, or when the effective sample
size is small, this assumption may be violated.

\subsection{4.2 TMLE-CF standard error inflation and
stability}\label{tmle-cf-standard-error-inflation-and-stability}

Cross-fitted TMLE (TMLE-CF) may exhibit larger or more variable standard
error estimates than standard TMLE in certain scenarios. Plausible
explanations include:

\begin{itemize}
\item
  \textbf{Fold instability}: When the sample is split into folds, some
  folds may contain fewer observations in regions of the covariate space
  where positivity is marginal. The resulting nuisance estimates can
  vary substantially across folds, inflating the variance of the
  cross-fitted estimator.
\item
  \textbf{Limited Super Learner library}: If the candidate learner
  library is insufficiently diverse, the Super Learner's performance may
  degrade when trained on smaller fold-specific training sets, leading
  to higher-variance nuisance estimates.
\item
  \textbf{Variance inflation under small effective sample size}:
  Cross-fitting can amplify the impact of near-positivity violations
  because each fold has a smaller effective sample size for estimating
  nuisance parameters in thin covariate regions.
\end{itemize}

These phenomena do not indicate a defect in cross-fitting per
se---cross-fitting provides stronger theoretical guarantees---but they
highlight the importance of adequate library specification and sample
size relative to the complexity of the nuisance functions. Useful
diagnostics include examining fold-specific propensity score
distributions, fold-specific ESS, and the variability of Super Learner
coefficients across folds.

\section{5. Discussion}\label{discussion}

\subsection{5.1 Summary of findings}\label{summary-of-findings}

\emph{{[}To be completed with simulation results.{]}}

\subsection{5.2 Outcome-blind simulation as a design validation
tool}\label{outcome-blind-simulation-as-a-design-validation-tool}

The outcome-blind simulation phase should be understood as a
\textbf{design validation} step within the clean-room governance
framework, not as a data-driven estimator selection procedure. Its role
is to confirm that the pre-specified analysis plan---estimator, Super
Learner library, truncation bounds, and diagnostic criteria---produces
estimates with acceptable bias, coverage, and variance calibration under
plausible data-generating conditions.

This distinction is important. If the outcome-blind phase were used to
select among estimators based on their simulated performance and then
apply the selected estimator to real outcome data, the procedure would
introduce a form of post-hoc analytic flexibility that the clean-room
framework is designed to prevent. Instead, the outcome-blind simulation
phase validates pre-specified choices and identifies potential problems
(e.g., poor coverage under informative censoring) that require plan
revisions before outcome access.

Pre-specified revisions triggered by simulation results---such as
expanding the Super Learner library or adjusting truncation bounds---are
documented in the decision log and are permissible under the clean-room
governance protocol precisely because they occur before outcome
unblinding.

\subsection{5.3 Limitations}\label{limitations}

Several limitations of the current simulation should be noted:

\begin{itemize}
\item
  \textbf{No competing risks.} The DGP does not model competing events
  such as death. In a real-world HCV-AKI study, death before AKI would
  be a competing risk that could bias cause-specific cumulative
  incidence estimates. Future extensions should incorporate competing
  risks, either by modeling them explicitly or by using subdistribution
  hazard approaches.
\item
  \textbf{Baseline confounding only.} The current DGP generates
  confounding through baseline covariates. Time-varying
  confounding---where covariates measured after treatment initiation
  affect both treatment continuation and outcome---is not simulated.
  Extending the simulation to include time-varying confounders would
  require longitudinal TMLE (LTMLE) estimation.
\item
  \textbf{Switching as censoring.} Treatment switching is handled by
  censoring at crossover, which approximates a hypothetical intervention
  that eliminates switching. This requires correct specification of the
  switching model for valid inference. An alternative would be to model
  switching as a time-varying treatment and use a sustained treatment
  strategy estimand, which would require a different estimation
  approach.
\item
  \textbf{Monte Carlo replicates.} The 200 replicates per scenario
  provide adequate power to detect large performance differences but may
  not resolve small differences in coverage or bias.
\end{itemize}

\subsection{5.4 Sensitivity analyses beyond the simulation
scenarios}\label{sensitivity-analyses-beyond-the-simulation-scenarios}

In a real-world application of this framework, the analysis plan would
include sensitivity analyses beyond those captured by the DGP scenarios.
Standard sensitivity analyses for targeted learning studies in
pharmacoepidemiology include:

\begin{itemize}
\item
  \textbf{Unmeasured confounding.} Because conditional exchangeability
  is untestable, sensitivity analyses quantify how strong an unmeasured
  confounder would need to be to explain away the observed effect.
  Approaches include the E-value {[}@VanderWeele2017{]}, which provides
  a minimum strength of association (on the risk ratio scale) that an
  unmeasured confounder would need to have with both treatment and
  outcome to fully explain the point estimate, and formal bias functions
  that parameterize the bias as a function of confounder prevalence and
  effect size.
\item
  \textbf{Truncation sensitivity.} The choice of truncation bounds for
  the propensity score (e.g., \([0.01, 0.99]\) vs.~\([0.05, 0.95]\)
  vs.~\([0.10, 0.90]\)) trades off bias against variance. Reporting
  results across multiple truncation levels demonstrates whether
  conclusions are robust to this analytic decision.
\item
  \textbf{Alternative Super Learner libraries.} Results may depend on
  the composition of the Super Learner library. Sensitivity to library
  specification can be assessed by running the analysis with alternative
  libraries (e.g., parametric only, non-parametric only, with or without
  HAL).
\item
  \textbf{Negative control outcomes and exposures.} A negative control
  outcome is one that should not be causally affected by the treatment;
  a negative control exposure is one that should not causally affect the
  outcome. Detecting associations with negative controls suggests
  residual confounding or systematic bias and serves as an empirical
  falsification test {[}@Lipsitch2010{]}.
\end{itemize}

\section{6. Conclusion}\label{conclusion}

\emph{{[}To be completed with simulation results. The conclusion should
summarize: (1) the performance of TMLE with Super Learner relative to
comparators; (2) the value of outcome-blind simulation as a design
validation tool; (3) practical recommendations for implementing
clean-room targeted learning in pharmacoepidemiologic safety
studies.{]}}

\section{7. Reproducibility}\label{reproducibility}

\subsection{7.1 Software environment}\label{software-environment}

All analyses are implemented in R. The following packages are required:

\begin{longtable}[]{@{}ll@{}}
\toprule\noalign{}
Package & Purpose \\
\midrule\noalign{}
\endhead
\bottomrule\noalign{}
\endlastfoot
\texttt{tmle} & Point-treatment TMLE estimation \\
\texttt{ltmle} & Longitudinal TMLE for survival outcomes \\
\texttt{lmtp} & Modified treatment policy estimation (alternative
TMLE) \\
\texttt{SuperLearner} & Ensemble machine learning for nuisance
parameters \\
\texttt{ranger} & Random forest learner for Super Learner library \\
\texttt{xgboost} & Gradient boosted trees for Super Learner library \\
\texttt{mgcv} & GAM learner for Super Learner library \\
\texttt{glmnet} & Penalized regression learner \\
\texttt{dbarts} & BART learner for Super Learner library \\
\texttt{hal9001} & Highly adaptive lasso (optional) \\
\texttt{tidyverse} & Data manipulation and visualization \\
\texttt{data.table} & Fast data operations for simulation loops \\
\texttt{here} & Project-relative file paths \\
\texttt{survival} & Cox PH regression comparator \\
\end{longtable}

\subsection{7.2 Configuration}\label{configuration}

The Super Learner library, truncation bounds, number of cross-validation
folds, and all DGP parameters are specified in the configuration files
archived with the repository. The appendix reproduces the DGP function
call and Super Learner specification used for each scenario.

Truncation bounds for the propensity score:
\([\delta_{\min}, \delta_{\max}] = [0.01, 0.99]\) (primary analysis),
with sensitivity analyses at \([0.05, 0.95]\) and \([0.10, 0.90]\).

\subsection{7.3 Seeds and
reproducibility}\label{seeds-and-reproducibility}

All simulation seeds are fixed and stored. The DGP function accepts a
\texttt{seed} argument that sets the random number generator state
before each replicate. Session information (R version, package versions,
platform) is captured and archived with each run.

\subsection{7.4 Unit tests}\label{unit-tests}

The test suite (\texttt{tests/test\_sim.R}) covers:

\begin{itemize}
\tightlist
\item
  \textbf{DGP validity}: verifies that generated datasets have expected
  dimensions, covariate distributions within specified tolerances,
  treatment prevalence near the target, and no missing values in
  required fields after cohort exclusions.
\item
  \textbf{Estimator smoke tests}: confirms that the TMLE, AIPW, and Cox
  pipelines run to completion on a small simulated dataset and produce
  estimates within the parameter space.
\item
  \textbf{Reproducibility}: confirms that running the DGP with the same
  seed produces identical datasets.
\end{itemize}

\subsection{7.5 Code availability}\label{code-availability}

All simulation code, DGP functions, estimation pipelines, and analysis
scripts are archived in the project repository with version control. The
repository structure follows the clean-room governance framework, with
separate directories for the DGP (\texttt{R/dgp/}), estimation pipeline
(\texttt{R/tmle/}), simulation driver (\texttt{simulation/}), and
reports (\texttt{reports/}).

\section{References}\label{references}

\phantomsection\label{refs}
\begin{CSLReferences}{0}{1}
\end{CSLReferences}

Berger ML, Sox H, Willke RJ, et al.~Good practices for real-world data
studies of treatment and/or comparative effectiveness: recommendations
from the joint ISPOR-ISPE Special Task Force on Real-World Evidence in
Health Care Decision Making. \emph{Pharmacoepidemiology and Drug
Safety}. 2017;26(9):1033-1039.

Chernozhukov V, Chetverikov D, Demirer M, et al.~Double/debiased machine
learning for treatment and structural parameters. \emph{The Econometrics
Journal}. 2018;21(1):C1-C68.

Hernan MA, Robins JM. Using big data to emulate a target trial when a
randomized trial is not available. \emph{American Journal of
Epidemiology}. 2016;183(8):758-764.

Lash TL, VanderWeele TJ, Haneuse S, Rothman KJ. \emph{Modern
Epidemiology}. 4th ed.~Wolters Kluwer; 2021.

Lipsitch M, Tchetgen Tchetgen E, Cohen T. Negative controls: a tool for
detecting confounding and bias in observational studies.
\emph{Epidemiology}. 2010;21(3):383-388.

Nosek BA, Ebersole CR, DeHaven AC, Mellor DT. The preregistration
revolution. \emph{Proceedings of the National Academy of Sciences}.
2018;115(11):2600-2606.

Petersen ML, van der Laan MJ. Causal models and learning from data:
integrating causal modeling and statistical estimation.
\emph{Epidemiology}. 2014;25(3):418-426.

Polley EC, van der Laan MJ. Super Learner in prediction. U.C. Berkeley
Division of Biostatistics Working Paper Series. 2010; Working Paper 266.

van der Laan MJ, Polley EC, Hubbard AE. Super Learner. \emph{Statistical
Applications in Genetics and Molecular Biology}. 2007;6(1):Article 25.

van der Laan MJ, Rubin D. Targeted maximum likelihood learning.
\emph{The International Journal of Biostatistics}. 2006;2(1):Article 11.

van der Laan MJ, Rose S. \emph{Targeted Learning: Causal Inference for
Observational and Experimental Data}. New York: Springer; 2011.

van der Laan MJ, Rose S. \emph{Targeted Learning in Data Science: Causal
Inference for Complex Longitudinal Studies}. New York: Springer; 2018.

VanderWeele TJ, Ding P. Sensitivity analysis in observational research:
introducing the E-value. \emph{Annals of Internal Medicine}.
2017;167(4):268-274.

Zheng W, van der Laan MJ. Cross-validated targeted minimum-loss-based
estimation. In: van der Laan MJ, Rose S, eds.~\emph{Targeted Learning}.
Springer; 2011:459-474.

\section{Appendix: DGP Configuration}\label{appendix-dgp-configuration}

\subsection{A.1 Simple scenario}\label{a.1-simple-scenario}

\begin{Shaded}
\begin{Highlighting}[]
\FunctionTok{generate\_hcv\_data}\NormalTok{(}
  \AttributeTok{N           =} \DecValTok{50000}\NormalTok{,}
  \AttributeTok{p\_sof       =} \FloatTok{0.36}\NormalTok{,}
  \AttributeTok{h0          =} \FloatTok{3e{-}4}\NormalTok{,}
  \AttributeTok{HR\_early    =} \FloatTok{1.25}\NormalTok{,}
  \AttributeTok{HR\_late     =} \FloatTok{0.70}\NormalTok{,}
  \AttributeTok{tau         =} \DecValTok{45}\NormalTok{,}
  \AttributeTok{max\_follow  =} \DecValTok{720}\NormalTok{,}
  \AttributeTok{risk\_window =} \DecValTok{30}\NormalTok{,}
  \AttributeTok{np\_hazard   =} \ConstantTok{FALSE}\NormalTok{,}
  \AttributeTok{dep\_censor  =} \ConstantTok{FALSE}\NormalTok{,}
  \AttributeTok{complexity  =} \ConstantTok{FALSE}\NormalTok{,}
  \AttributeTok{switch\_on   =} \ConstantTok{FALSE}\NormalTok{,}
  \AttributeTok{seed        =} \SpecialCharTok{\textless{}}\NormalTok{replicate}\SpecialCharTok{{-}}\NormalTok{specific}\SpecialCharTok{\textgreater{}}
\NormalTok{)}
\end{Highlighting}
\end{Shaded}

\subsection{A.2 Complex scenario (Scenario 4: treatment
switching)}\label{a.2-complex-scenario-scenario-4-treatment-switching}

\begin{Shaded}
\begin{Highlighting}[]
\FunctionTok{generate\_hcv\_data}\NormalTok{(}
  \AttributeTok{N           =} \DecValTok{50000}\NormalTok{,}
  \AttributeTok{p\_sof       =} \FloatTok{0.36}\NormalTok{,}
  \AttributeTok{h0          =} \FloatTok{5e{-}5}\NormalTok{,}
  \AttributeTok{HR\_early    =} \FloatTok{1.25}\NormalTok{,}
  \AttributeTok{HR\_late     =} \FloatTok{0.70}\NormalTok{,}
  \AttributeTok{tau         =} \DecValTok{45}\NormalTok{,}
  \AttributeTok{max\_follow  =} \DecValTok{720}\NormalTok{,}
  \AttributeTok{risk\_window =} \DecValTok{30}\NormalTok{,}
  \AttributeTok{np\_hazard   =} \ConstantTok{TRUE}\NormalTok{,}
  \AttributeTok{dep\_censor  =} \ConstantTok{TRUE}\NormalTok{,}
  \AttributeTok{complexity  =} \ConstantTok{TRUE}\NormalTok{,}
  \AttributeTok{switch\_on   =} \ConstantTok{TRUE}\NormalTok{,}
  \AttributeTok{lambda\_sw0  =} \FloatTok{2.5e{-}5}\NormalTok{,}
  \AttributeTok{gamma\_A     =} \FloatTok{0.60}\NormalTok{,}
  \AttributeTok{gamma\_ckd   =} \FloatTok{0.40}\NormalTok{,}
  \AttributeTok{seed        =} \SpecialCharTok{\textless{}}\NormalTok{replicate}\SpecialCharTok{{-}}\NormalTok{specific}\SpecialCharTok{\textgreater{}}
\NormalTok{)}
\end{Highlighting}
\end{Shaded}

\subsection{A.3 Super Learner library
specification}\label{a.3-super-learner-library-specification}

\begin{Shaded}
\begin{Highlighting}[]
\NormalTok{SL.library }\OtherTok{\textless{}{-}} \FunctionTok{c}\NormalTok{(}
  \StringTok{"SL.glm"}\NormalTok{,}
  \StringTok{"SL.gam"}\NormalTok{,}
  \StringTok{"SL.ranger"}\NormalTok{,}
  \StringTok{"SL.xgboost"}\NormalTok{,}
  \StringTok{"SL.dbarts"}\NormalTok{,}
  \StringTok{"SL.mean"}
\NormalTok{)}

\CommentTok{\# Cross{-}validation folds for Super Learner}
\NormalTok{cvControl }\OtherTok{\textless{}{-}} \FunctionTok{list}\NormalTok{(}\AttributeTok{V =} \DecValTok{10}\NormalTok{)}

\CommentTok{\# Truncation bounds for propensity scores}
\NormalTok{g.bounds }\OtherTok{\textless{}{-}} \FunctionTok{c}\NormalTok{(}\FloatTok{0.01}\NormalTok{, }\FloatTok{0.99}\NormalTok{)}
\end{Highlighting}
\end{Shaded}





\end{document}
